\section{Methodology}

\subsection{System Overview}
The battery thermal digital twin pipeline is composed of five major components: dataset generation via Simulink, incremental preprocessing and scaling, baseline Multi-Layer Perceptron (MLP) training, Physics-Informed Neural Network (PINN) training, and autoregressive rollout evaluation. 

The dataset is generated from a high-fidelity Simulink electro-thermal model simulating cell behavior under dynamic loads. Preprocessing calculates the temperature change $\Delta T_t = T_{t+1} - T_t$ and time step $dt_t = t_{t+1} - t_t$ independently for each drive cycle file. This prevents artificial transitions across cycle boundaries. Features are scaled using standard normalization fitted incrementally to avoid memory bottlenecks. The models take current states as inputs and predict the temperature delta over the time step.

\begin{figure}[htbp]
\centering
\begin{tikzpicture}[
    auto,
    block/.style={rectangle, draw, fill=blue!10, text width=2.4cm, text centered, rounded corners, minimum height=0.9cm, font=\scriptsize},
    line/.style={draw, -latex, thick}
]
    \node [block] (simulink) {Simulink Data\\Generation};
    \node [block, below=0.4cm of simulink] (preprocess) {Preprocessing \&\\Scaling};
    \node [block, below=0.4cm of preprocess] (train) {MLP / PINN\\Training};
    \node [block, below=0.4cm of train] (eval) {Autoregressive\\Rollout};
    
    \path [line] (simulink) -- (preprocess);
    \path [line] (preprocess) -- (train);
    \path [line] (train) -- (eval);
\end{tikzpicture}
\caption{The battery thermal digital twin pipeline from Simulink data generation to autoregressive rollout evaluation.}
\label{fig:pipeline}
\end{figure}

\subsection{Lumped Battery Thermal Model}
The physics constraints embedded within the PINN are derived from the governing ordinary differential equation (ODE) of a lumped-parameter thermal model \cite{bernardi1985general}. The cell energy balance is formulated in \eqref{eq:governing} as:
\begin{equation}
m C_p \frac{dT}{dt} = P_{\text{loss}} - Q_{\text{loss}}
\label{eq:governing}
\end{equation}
where $m$ is the mass of the battery cell (kg), $C_p$ is the specific heat capacity (J/kg/K), $T$ is the cell temperature (K), $P_{\text{loss}}$ is the heat generation rate (W) due to internal Joule heating and polarization losses, and $Q_{\text{loss}}$ is the heat transfer rate to the ambient environment (W).

Assuming Newton's law of cooling for convection, the heat loss $Q_{\text{loss}}$ is expressed in \eqref{eq:convection} as:
\begin{equation}
Q_{\text{loss}} = hA (T - T_{\text{amb}})
\label{eq:convection}
\end{equation}
where $hA$ is the convective heat transfer coefficient multiplied by the cell surface area (W/K), and $T_{\text{amb}}$ is the ambient temperature (K). 

To account for cell scaling variations across different experiments in the dataset, we incorporate a dimensionless scaling factor, $\text{MassScale}$. Combining \eqref{eq:governing} and \eqref{eq:convection}, the physical rate of temperature change is modeled in \eqref{eq:rate_of_change} as:
\begin{equation}
\frac{dT}{dt} = \frac{P_{\text{loss}} - hA (T - T_{\text{amb}})}{m C_p \cdot \text{MassScale}}
\label{eq:rate_of_change}
\end{equation}

In this study, the base thermal capacity is set to $m C_p = 900.0$ J/K, which matches the physical properties of standard lithium-ion cells used in the Simulink generator.

\begin{figure}[htbp]
\centering
\begin{tikzpicture}[
    node distance=1.2cm,
    node/.style={circle, draw, fill=orange!25, minimum size=0.8cm, font=\scriptsize},
    ambient/.style={circle, draw, fill=blue!10, minimum size=0.8cm, font=\scriptsize},
    resistor/.style={rectangle, draw, fill=gray!15, minimum height=0.3cm, minimum width=0.7cm, font=\tiny}
]
    \node[node] (core) {$T_c$};
    \node[resistor, right=0.6cm of core] (Rcs) {$R_{cs}$};
    \node[node, right=0.6cm of Rcs] (case) {$T_s$};
    \node[resistor, right=0.6cm of case] (Rconv) {$R_{\text{conv}}$};
    \node[ambient, right=0.6cm of Rconv] (amb) {$T_{\text{amb}}$};
    
    \draw[thick] (core) -- (Rcs);
    \draw[thick, ->] (Rcs) -- (case);
    \draw[thick] (case) -- (Rconv);
    \draw[thick, ->] (Rconv) -- (amb);
    
    \node[below=0.1cm of core, font=\tiny, text centered] {Core Heat $P_{\text{loss}}$};
    \node[below=0.1cm of case, font=\tiny, text centered] {Casing};
\end{tikzpicture}
\caption{Equivalent thermal circuit showing the dual-node core-to-casing heat flow in the Simulink environment.}
\label{fig:thermal_circuit}
\end{figure}

\subsection{Baseline Multi-Layer Perceptron (MLP)}
The baseline neural network is a fully connected Multi-Layer Perceptron designed to map the battery electro-thermal features to the temperature change. The model inputs are represented by the feature vector in \eqref{eq:features}:
\begin{equation}
\mathbf{x}_t = [T_t, I_t, V_t, \text{SOC}_t, P_{\text{loss}, t}, T_{\text{amb}, t}, \text{MassScale}_t]^T
\label{eq:features}
\end{equation}
where $I_t$ is current, $V_t$ is voltage, and $\text{SOC}_t$ is state-of-charge. The target output is the physical temperature delta in \eqref{eq:target}:
\begin{equation}
y_t = \Delta T_t = T_{t+1} - T_t
\label{eq:target}
\end{equation}

The network architecture is defined by the following layer structure:
\begin{align*}
\mathbf{h}_1 &= \text{ReLU}(\mathbf{W}_1 \mathbf{x}_t + \mathbf{b}_1) \quad (\mathbf{W}_1 \in \mathbb{R}^{64 \times 7}) \\
\mathbf{h}_2 &= \text{ReLU}(\mathbf{W}_2 \mathbf{h}_1 + \mathbf{b}_2) \quad (\mathbf{W}_2 \in \mathbb{R}^{64 \times 64}) \\
\mathbf{h}_3 &= \text{ReLU}(\mathbf{W}_3 \mathbf{h}_2 + \mathbf{b}_3) \quad (\mathbf{W}_3 \in \mathbb{R}^{32 \times 64}) \\
\Delta \hat{T}_t &= \mathbf{W}_4 \mathbf{h}_3 + \mathbf{b}_4 \quad (\mathbf{W}_4 \in \mathbb{R}^{1 \times 32})
\end{align*}
The network is optimized using the Mean Squared Error (MSE) data loss in \eqref{eq:loss_data}:
\begin{equation}
\mathcal{L}_{\text{data}} = \frac{1}{N} \sum_{i=1}^N \left( \Delta T_i - \Delta \hat{T}_i \right)^2
\label{eq:loss_data}
\end{equation}

\subsection{Physics-Informed Neural Network (PINN) Formulation}
The PINN shares the architectural design of the baseline MLP but incorporates the governing physical ODE equation \eqref{eq:rate_of_change} as a soft regularization term in its objective function. This regularizes the neural network search space toward thermodynamically plausible solutions via a composite loss function in \eqref{eq:loss_total}:
\begin{equation}
\mathcal{L}_{\text{total}} = \mathcal{L}_{\text{data}} + \lambda_{\text{phys}} \mathcal{L}_{\text{physics}}
\label{eq:loss_total}
\end{equation}
where $\lambda_{\text{phys}}$ is a regularization weight, set to $0.01$ in our training configuration.

The physics loss term measures the residual of the lumped single-node thermal ODE over each mini-batch. To evaluate the residual, the predicted temperature delta in scaled units is inverted back to physical units via \eqref{eq:unscale}:
\begin{equation}
\Delta \hat{T}_i^{\text{phys}} = \Delta \hat{T}_i \cdot \sigma_y + \mu_y
\label{eq:unscale}
\end{equation}
where $\mu_y$ and $\sigma_y$ are the empirical target mean and standard deviation.

The physics residual $R_i$ for sample $i$ is calculated in physical units (K/s) in \eqref{eq:residual} as:
\begin{equation}
R_i = \frac{\Delta \hat{T}_i^{\text{phys}}}{dt_i} - \frac{P_{\text{loss}, i} - hA (T_i - T_{\text{amb}, i})}{m C_p \cdot \text{MassScale}_i}
\label{eq:residual}
\end{equation}
where $T_i$, $P_{\text{loss}, i}$, $T_{\text{amb}, i}$, and $\text{MassScale}_i$ are unscaled inputs, and $dt_i$ is the time step size in seconds. The physics regularization loss is computed in \eqref{eq:loss_physics} as:
\begin{equation}
\mathcal{L}_{\text{physics}} = \frac{1}{N} \sum_{i=1}^N R_i^2
\label{eq:loss_physics}
\end{equation}

\subsection{Trainable Parameter Parameterization}
The convective heat transfer coefficient, $hA$, is treated as a trainable model parameter. Because a negative value of $hA$ would represent non-physical environmental heating when $T > T_{\text{amb}}$, the parameter is constrained to positive values. Simple parameter clamping halts gradient flow when hitting the boundary ($hA \to 0$). To maintain continuous gradient updates, we parameterize $hA$ using a raw real-valued variable $hA_{\text{raw}}$ passed through a softplus activation in \eqref{eq:softplus}:
\begin{equation}
hA = \text{softplus}(hA_{\text{raw}}) = \ln(1 + e^{hA_{\text{raw}}})
\label{eq:softplus}
\end{equation}

To initialize $hA$ at a baseline value of $hA_{\text{init}} = 0.1$ W/K, $hA_{\text{raw}}$ is initialized via \eqref{eq:raw_init} as:
\begin{equation}
hA_{\text{raw, init}} = \ln(e^{hA_{\text{init}}} - 1)
\label{eq:raw_init}
\end{equation}
This formulation guarantees that $hA > 0$ throughout training. Because the training targets reflect cell surface temperature from a dual-node Simulink simulation, the learned parameter $hA$ represents an \textbf{effective convective heat transfer parameter} within the single-node ODE approximation.
